% !TeX root = ../main.tex

\chapter{Complex Extension}

% =================================================

Another possible generalisation is to allow

\[
\lambda\in\mathbb{C}.
\]

If

\[
\lambda=i\theta,
\]

then

\[
\phi_n(x)
=
e^{i\theta x^n}.
\]

Since

\[
|e^{i\theta x^n}|=1,
\]

the blocks no longer grow or decay.

Instead they oscillate.

Thus the same structure becomes

\[
\boxed{
F(x)
=
\sum_{n\geq1}
c_ne^{i\theta x^n}.
}
\]

Now the problem is no longer dominated by exponential growth but by
cancellation and oscillation.

This creates a possible connection with polynomial-phase exponential
functions and harmonic analysis.

% =================================================
\section{What is Inherited and What May Be Interesting}
% =================================================

The identity

\[
\phi_n(x^m)=\phi_{nm}(x)
\]

is ultimately inherited from

\[
(x^m)^n=x^{mn}.
\]

Therefore this identity by itself should not be considered a new
phenomenon.

However, several richer structures appear after the blocks are placed
inside linear combinations.

For example,

\[
F(x)
=
\sum c_n\phi_n(x)
\]

allows the composition operators \(T_m\) to transform entire
coefficient sequences.

Combining these operators produces Dirichlet convolution.

Taylor expansion produces divisor sums independently.

The centered family

\[
e^{\lambda x^n}-1
\]

also gives a triangular transformation of ordinary power series.

These interactions appear to be more mathematically substantial than
the original identity alone.

% =================================================
